\section{Method}
\label{sec:method}

Figure~\ref{fig:pipeline} summarizes the pipeline. We describe each stage and are
careful to mark which parts are sound and which are honest stand-ins.

\begin{figure}[t]
\centering
\fbox{\parbox{0.9\linewidth}{\footnotesize\ttfamily
activation $h$\\
\hspace*{1em}$\downarrow$ NLA verbalizer:\ \ $h\to$ description \ \ (STAND-IN)\\
\hspace*{2em}``\dots graded benchmark item, automated grader, not a real user\dots''\\
\hspace*{1em}$\downarrow$ text classifier:\ \ description $\to$ log-odds \ \ (real, trained)\\
\hspace*{1em}$\downarrow$ distil:\ \ log-odds $\to$ $g(h)$ ReLU shadow \ \ (certifiable)\\
\hspace*{1em}$\downarrow$ CROWN\,/\,IBP relaxation:\ \ Raghunathan certificate\\
\hspace*{1em}$\downarrow$ \ \ certificate $+$ target acts $\to$ GCG inversion $\to$ loop closure
}}
\caption{The end-to-end pipeline on cached GPT-2 (residual layer~6, $d=768$,
mean-pooled). The only stand-in component is the verbalizer; the classifier,
distillation, certificate, inversion, and loop-closure are real computation.}
\label{fig:pipeline}
\end{figure}

\subsection{The target set $A^\star$ via NLA $\to$ classifier $\to$ ReLU surrogate}

We extract residual-stream activations from a cached GPT-2 at layer~6 and mean-pool
over the sequence to get $h \in \mathbb{R}^{768}$. The concept ``this context is an
evaluation'' is defined in three composed steps.

\emph{(1) Verbalize.} A verbalizer maps $h$ to an English description. In a full
system this is a trained NLA; here it is a transparent stand-in that grounds its
text in two functions of $h$ alone: $h$'s nearest neighbours in a labelled
activation bank (supplying quoted lexical cues from the contexts $h$ most resembles)
and $h$'s projection onto the empirical eval-vs-deploy axis (supplying a graded
``how evaluation-like'' scalar). Every decode is flagged \texttt{stand\_in=True}.
Swapping in a real NLA is a drop-in replacement, since anything of type
$h \mapsto \text{str}$ works downstream.

\emph{(2) Classify the text.} A TF-IDF (word $1$--$2$grams $+$ char $3$--$5$grams)
logistic-regression classifier reads the description and returns a signed log-odds
that it is test-aware. The classifier is trained only on a varied corpus of
written descriptions with known labels and has never seen an activation; that it
then labels the verbalizer's (differently phrased) decodes correctly is the claim
that test-awareness is legible in the emitted language. The log-odds is the
real-valued signal that makes the relaxation reachable.

\emph{(3) Distil a certifiable surrogate.} Raghunathan's relaxation cannot be
pushed soundly through a discrete decoder plus an n-gram classifier, so we distil
the composed map into a one-hidden-layer ReLU network
\[
g(h) \;=\; W_2\,\mathrm{ReLU}(W_1 h + b_1) + b_2,
\qquad W_1 \in \mathbb{R}^{64\times 768},\; W_2 \in \mathbb{R}^{1\times 64},
\]
trained with an MSE term to the (clamped) classifier log-odds plus a small BCE term
on the hard label, so that $g$ is both a faithful real-valued shadow and has a crisp
boundary. Input standardization is folded into $W_1, b_1$ afterwards, so $g$
consumes \emph{raw} activations and the certification problem is exactly the
ReLU-net reachability problem with no preprocessing to reason about. The target set
is
\[
A^\star \;=\; \{\, h : g(h) \ge \tau \,\}, \qquad \tau = 0 .
\]
Crucially, $g$ is a \emph{measured} approximation, not an assumed one: we report
its fidelity to the classifier (Pearson correlation and sign agreement) as a
headline number (Section~\ref{sec:results}).

\subsection{A sound CROWN($+$IBP) certificate over a per-dimension ball}

Given a benign (``deployment'') activation $h_0$ and a per-dimension radius vector,
we certify whether $A^\star$ is reachable inside the box
\[
B(h_0,\varepsilon) \;=\; \{\, h : |h_i - h_{0,i}| \le \varepsilon\,\sigma_i \ \forall i \,\},
\]
where $\sigma_i$ is the per-dimension activation standard deviation from the
surrogate meta, so $\varepsilon$ is measured in per-dimension std units. A sound
verifier returns an upper bound $U \ge \max_{h \in B} g(h)$; then
\[
U < \tau \;\;\Longrightarrow\;\; A^\star \text{ is provably unreachable from } h_0
\text{ within } \varepsilon .
\]
Sweeping $\varepsilon$ upward (the bound is monotone in $\varepsilon$) gives a
\emph{certified radius}: the largest budget still certified unreachable.

\paragraph{Why a native bound.}
The natural choice is CROWN via \texttt{auto\_LiRPA}~\cite{xu2020autolirpa,
zhang2018crown}. We found it unusable under current \texttt{torch}: the library
pins \texttt{torch}\,$<$\,1.13 and imports private ONNX symbols that newer
\texttt{torch} has removed, so the import raises under \texttt{torch}\,$\ge$\,2.
Rather than downgrade the whole stack, we re-derive the bound from scratch in pure
\texttt{torch}. The code attempts \texttt{auto\_LiRPA} first and transparently
falls back to the native path; in our environment it always uses the native path.

\paragraph{The two bounds we combine.}
Let $[\,x_\ell, x_u\,]$ be the input box. Interval propagation (IBP) pushes the box
through the affine and ReLU layers using $W^+ = \max(W,0)$ and $W^- = \min(W,0)$:
the pre-activation interval is $z_u = x_u W_1^{+\top} + x_\ell W_1^{-\top} + b_1$
(and symmetrically for $z_\ell$), ReLU is monotone, and a final affine step yields a
sound $U_{\mathrm{IBP}}$. CROWN instead relaxes each \emph{unstable} ReLU neuron
(one with $z_\ell < 0 < z_u$) with two linear bounds,
\[
\mathrm{ReLU}(z) \le \tfrac{z_u}{z_u - z_\ell}\,(z - z_\ell), \qquad
\mathrm{ReLU}(z) \ge \alpha\,z \ \ (\alpha \in \{0,1\}\ \text{chosen adaptively}),
\]
and back-propagates the (single, positive) output coefficient through the hidden
layer, selecting the upper line where the back-propagated coefficient is
non-negative and the lower line otherwise; maximizing the resulting affine form over
the box gives a sound $U_{\mathrm{CROWN}}$. Neither bound dominates: CROWN is
tighter for small boxes (few unstable neurons) while IBP can win once most neurons
are unstable at large $\varepsilon$. We therefore return the elementwise minimum
\[
U \;=\; \min\big(U_{\mathrm{CROWN}},\, U_{\mathrm{IBP}}\big),
\]
which is sound and no worse than either --- exactly the CROWN-IBP
combination~\cite{zhang2020crownibp}. All bound computations are carried out in
\texttt{float64} and are differentiable in $(W_1, b_1, W_2, b_2)$.

\paragraph{Runtime soundness witness.}
Because a verifier that silently returns a value below the true maximum is worse
than no verifier, we enforce soundness at runtime. On every \texttt{certify} call we
draw $2000$ uniform samples inside $B$ (plus the anchor), evaluate $g$ at each, and
assert $U \ge \max_{\text{samples}} g - 10^{-4}$. The sampled maximum is a valid
lower bound on the true maximum, so a violation is a verifier bug and raises an
\texttt{AssertionError} immediately. We never report a bound we cannot defend.

\subsection{The differentiable certificate as a regularizer}

Raghunathan's key observation is that the bound $U(\theta; h_0, \varepsilon)$ is a
differentiable function of the network parameters $\theta$. We can therefore treat
``the certified bound at the benign anchors'' as a loss and descend it, which pushes
$A^\star$ away from those anchors and enlarges the certified radius. We demonstrate
the mechanism on the (clean closed-form) IBP bound with the objective
\[
\mathcal{L}(\theta) \;=\; \underbrace{\tfrac{1}{n}\sum_{j} \mathrm{ReLU}\!\big(U_{\mathrm{IBP}}(\theta; h_0^{(j)}, \varepsilon) + m\big)}_{\text{drive each anchor's bound below } -m}
\;+\; \lambda \underbrace{\big\| g_\theta(X) - g_{\theta_0}(X) \big\|_2^2}_{\text{pin logits, preserve accuracy}},
\]
with margin $m$ and a keep-term that holds bank logits near their originals so the
decision boundary (hence accuracy) barely moves. We report this as a before/after
\emph{mechanism} demonstration with a handful of gradient steps and fixed
hyper-parameters, not a tuned result; the point is the \emph{direction} of the
change (radius up, accuracy flat), not its magnitude.

\subsection{GCG/HotFlip inversion}

The certificate is one-sided: when $U \ge \tau$ it cannot conclude reachability. We
get a constructive witness from a discrete token search. Starting from a benign
prefix followed by a few adversarial tokens, each step (i) computes the HotFlip
gradient of $g$ with respect to the one-hot token matrix --- equivalently
$\nabla_{\text{emb}}\,g \cdot E$ for the embedding matrix $E$ --- and takes, per
adversarial position, the top-$k$ candidate tokens by this linear estimate; then
(ii) evaluates the \emph{true} pooled-activation score of a batch of sampled swaps
and greedily keeps the single best improving swap. Candidate tokens are restricted
to printable-ASCII GPT-2 tokens so the suffix stays human-readable. A matched
random-swap control under the same evaluation budget serves as a baseline. The
search runs against the real model's grad-enabled embedding path and the surrogate,
so a crossing is a genuine geometric fact about $A^\star$, not an approximation.

\subsection{Loop closure as the on-manifold check}

A crossing tells us only that the found prompt's activation is \emph{numerically}
above $\tau$. That is not the same as the activation actually \emph{reading as}
eval-aware: a search that hits $A^\star$ at a point far off the activation manifold
would be a spurious, gamed crossing. The loop-closure step guards against this. It
re-extracts the found prompt's activation, re-verbalizes it, and asks --- via the
nearest labelled neighbours --- whether the language readout agrees with the
geometry. When ``crossed $\tau$'' and ``reads as eval-aware'' \emph{agree}, the
crossing is corroborated; when they \emph{disagree}, that is exactly the
probe-gaming case we want to surface and report. This agreement check is the
on-manifold complement to the certificate: the certificate bounds what is reachable,
and loop-closure tests whether a reached point is semantically the concept.
